%!TEX root = main.tex
% =============================================================================
%  preamble.tex —— 2026 CUMCM C 题论文模板导言区
%  格式依据：《全国大学生数学建模竞赛论文格式规范（2026 年修订稿）》format2026.doc
%   第一条 A4、页边距≥2.5cm、左侧装订        -> geometry
%   第三条 摘要≤1页、页码自摘要页起、页脚中部、阿拉伯数字连续 -> fancyhdr
%   第四条 正文≤30页、不要目录                -> 不加载 \tableofcontents
%   第六条 摘要页/正文/附录不得出现身份信息    -> 不设 \author，PDF 元数据留空
%   第七条 引用需列参考文献并在正文标注        -> thebibliography（GB/T 7714）
% =============================================================================

% ---------------------------------------------------------------------------
% 1. 页面与版式  —— 第一条：A4，上下左右各留 ≥2.5cm，左侧装订
% ---------------------------------------------------------------------------
\usepackage[
  a4paper,
  top=2.5cm, bottom=2.5cm, left=2.5cm, right=2.5cm,
  footskip=1.0cm      % 页脚落在下边距内，正文块本身仍保有 2.5cm 边距
]{geometry}

% 行距：正文约 1.2 倍（格式规范第八条不对字号、行距作统一要求）。
% 正文有 30 页硬上限，此处取值已为「表格与图填入真实结果」留出余量；
% 若确实需要更松的行距，请先确认正文仍不超过 30 页。
\linespread{1.0}
% 允许 TeX 在极窄的表格列中略微拉伸以消除溢出（不影响观感）
\setlength{\emergencystretch}{2em}
% 段间距为 0（中文论文用首行缩进分段）；浮动体与标题间距一并收紧，
% 以便正文容纳更多内容而不牺牲可读性
\setlength{\parskip}{0pt}
\setlength{\textfloatsep}{9pt plus 2pt minus 2pt}
\setlength{\intextsep}{9pt plus 2pt minus 2pt}
\setlength{\floatsep}{9pt plus 2pt minus 2pt}
\setlength{\abovecaptionskip}{4pt}
\setlength{\belowcaptionskip}{0pt}

% ---------------------------------------------------------------------------
% 2. 页码  —— 第三条：自摘要页起、页脚中部、阿拉伯数字从 1 连续编号
%             摘要页即电子版第 1 页（第十条：承诺书与编号页不放电子版）
% ---------------------------------------------------------------------------
\usepackage{fancyhdr}
\pagestyle{fancy}
\fancyhf{}
\fancyfoot[C]{\normalfont\thepage}
\renewcommand{\headrulewidth}{0pt}
\renewcommand{\footrulewidth}{0pt}
\fancypagestyle{plain}{%
  \fancyhf{}\fancyfoot[C]{\normalfont\thepage}%
  \renewcommand{\headrulewidth}{0pt}\renewcommand{\footrulewidth}{0pt}%
}
\pagenumbering{arabic}

% ---------------------------------------------------------------------------
% 3. 数学、表格、图形
% ---------------------------------------------------------------------------
\usepackage{amsmath,amssymb,amsthm,mathrsfs,bm}
\allowdisplaybreaks[2]

% booktabs 的横线间距一并收紧，减少表格占用的纵向空间
\usepackage{booktabs,longtable,tabularx,array,multirow,makecell}
\setlength{\tabcolsep}{4pt}
\setlength{\abovedisplayskip}{4pt plus 2pt minus 2pt}
\setlength{\belowdisplayskip}{4pt plus 2pt minus 2pt}
\setlength{\abovedisplayshortskip}{2pt plus 1pt}
\setlength{\belowdisplayshortskip}{2pt plus 1pt}
\setlength{\aboverulesep}{0.18ex}
\setlength{\belowrulesep}{0.18ex}
\newcolumntype{C}[1]{>{\centering\arraybackslash}p{#1}}
\newcolumntype{L}[1]{>{\raggedright\arraybackslash}p{#1}}
\newcolumntype{Y}{>{\centering\arraybackslash}X}
\renewcommand{\arraystretch}{0.95}
% 浮动体（图/表）排布参数：本文图表较多，放宽参数可显著减少"半空页"
\usepackage{graphicx}
\setcounter{topnumber}{5}
\setcounter{bottomnumber}{3}
\setcounter{totalnumber}{8}
\renewcommand{\topfraction}{0.92}
\renewcommand{\bottomfraction}{0.85}
\renewcommand{\textfraction}{0.03}
\renewcommand{\floatpagefraction}{0.85}
\renewcommand{\dbltopfraction}{0.92}

\usepackage{subcaption}
\usepackage[font=small,labelsep=quad,skip=4pt]{caption}
\usepackage{enumitem}
\setlist{nosep,leftmargin=2.2em}
\usepackage{xcolor}
\usepackage{tikz}
\usetikzlibrary{arrows.meta,positioning,calc,patterns,shapes.geometric,fit,backgrounds,decorations.pathreplacing,decorations.markings}
\usepackage{pgfplots}
\pgfplotsset{compat=1.16}
\pgfplotsset{every axis/.append style={font=\small}}

% 算法
\usepackage{algorithm}
\usepackage{algpseudocode}
\renewcommand{\algorithmicrequire}{\textbf{输入：}}
\renewcommand{\algorithmicensure}{\textbf{输出：}}
\floatname{algorithm}{算法}

% ---------------------------------------------------------------------------
% 4. 定理类环境（中文名）
% ---------------------------------------------------------------------------
\theoremstyle{plain}
\newtheorem{proposition}{命题}
\newtheorem{lemma}{引理}
\newtheorem{corollary}{推论}
\theoremstyle{definition}
\newtheorem{definition}{定义}
\renewcommand{\proofname}{\heiti 证明}
\renewcommand{\qedsymbol}{$\square$}

% ---------------------------------------------------------------------------
% 5. 章节标题样式：一级「一、」三号黑体居中；二级 四号黑体；三级 小四黑体
% ---------------------------------------------------------------------------
\ctexset{
  section = {
    format     = \zihao{3}\heiti\centering,
    name       = {,、},
    number     = \chinese{section},
    aftername  = {},
    beforeskip = 1.2ex plus .4ex minus .2ex,
    afterskip  = 0.8ex plus .2ex,
  },
  subsection = {
    format     = \zihao{4}\heiti\raggedright,
    number     = \arabic{section}.\arabic{subsection},
    aftername  = \quad,
    beforeskip = 1.0ex plus .4ex minus .2ex,
    afterskip  = 0.5ex plus .2ex,
  },
  subsubsection = {
    format     = \zihao{-4}\heiti\raggedright,
    number     = \arabic{section}.\arabic{subsection}.\arabic{subsubsection},
    aftername  = \quad,
    beforeskip = 1.0ex plus .4ex minus .2ex,
    afterskip  = 0.5ex plus .2ex,
  },
  paragraph = {
    format     = \zihao{-4}\heiti\raggedright,
    aftername  = \quad,
  },
}

% ---------------------------------------------------------------------------
% 6. 源码清单（附录）—— 关键：basicstyle 中加入 \songti，
%    使代码中的中文注释用已嵌入的中文字体排版；否则中文会丢失或字体不嵌入。
% ---------------------------------------------------------------------------
\usepackage{listings}
\definecolor{cKeyword}{RGB}{0,80,160}
\definecolor{cComment}{RGB}{0,120,80}
\definecolor{cString} {RGB}{170,60,0}
\definecolor{cGray}   {RGB}{120,120,120}
\lstset{
  basicstyle       = \ttfamily\footnotesize\songti,
  keywordstyle     = \color{cKeyword}\bfseries,
  commentstyle     = \color{cComment}\itshape,
  stringstyle      = \color{cString},
  numberstyle      = \tiny\color{cGray},
  numbers          = left,
  numbersep        = 7pt,
  frame            = single,
  framerule        = 0.4pt,
  rulecolor        = \color{cGray},
  breaklines       = true,
  breakatwhitespace= false,
  columns          = flexible,
  keepspaces       = true,
  showstringspaces = false,
  tabsize          = 4,
  xleftmargin      = 2.2em,
  framexleftmargin = 2.2em,
  aboveskip        = 1.2ex,
  belowskip        = 1.2ex,
}
\lstdefinestyle{py}{language=Python}
\lstdefinestyle{plain}{language=}

% ---------------------------------------------------------------------------
% 7. 超链接与 PDF 元数据（第六条：元数据不得含身份信息）
% ---------------------------------------------------------------------------
% 黑体改用 PingFang SC：STHeiti 在 macOS 下 ToUnicode 生成失败会导致 PDF 文字无法复制
% 正文与黑体均启用伪粗体：macOS 的 Songti SC Bold 不嵌入 ToUnicode，
% 直接用 \textbf 会让加粗中文在复制时变成乱码；伪粗体由常规字重合成，可正常复制。
\setCJKmainfont{Songti SC}[AutoFakeBold=2.2]
\setCJKsansfont{Songti SC}[AutoFakeBold=2.2]
\setCJKmonofont{Songti SC}[AutoFakeBold=2.2]
\setCJKfamilyfont{zhhei}{PingFang SC}[BoldFont=PingFang SC Semibold]
\renewcommand{\heiti}{\CJKfamily{zhhei}}
\usepackage[hidelinks,bookmarks=true,bookmarksnumbered=false]{hyperref}
\hypersetup{
  pdfauthor  = {},
  pdfcreator = {},
  pdfsubject = {},
  pdftitle   = {微网与外部电网电力调控策略},
}

% ---------------------------------------------------------------------------
% 8. 写作期与定稿期开关
%    \drafttrue  ：显示【写作提示】与红色 {待计算} 占位符（写作期默认）
%    \draftfalse ：隐藏全部写作提示；占位符变为醒目红底，便于定稿前扫尾
%    两种进入定稿模式的方式（任选）：
%      (a) 把下面一行改成 \draftfalse
%      (b) 在论文目录下执行  touch final.flag   （便于反复切换与自动检查）
%    定稿前务必运行 tools/check_paper.py
% ---------------------------------------------------------------------------
\newif\ifdraft
\drafttrue
\IfFileExists{final.flag}{\draftfalse}{}

\definecolor{cGuide}{RGB}{90,90,90}
\definecolor{cNote} {RGB}{0,90,180}
\definecolor{cPh}   {RGB}{200,0,0}
\definecolor{cEvid} {RGB}{120,60,150}

\usepackage{environ}

% 带圈数字（①②③…）、罗马数字等符号在默认字体分类下会被交给西文字体而缺字，
% 这里把它们划入 CJK 字符类，交给中文字体排版（仅 XeLaTeX + xeCJK 生效）
\ifdefined\xeCJKsetcharclass
  \xeCJKsetcharclass{"2460}{"24FF}{1}   % ①②③…⒈⒉⒊…
  \xeCJKsetcharclass{"2160}{"217F}{1}   % Ⅰ Ⅱ Ⅲ … ⅰ ⅱ ⅲ
  \xeCJKsetcharclass{"3000}{"303F}{1}   % CJK 标点
\fi

% 【写作提示】块：draft 模式显示，定稿模式整块消失
\NewEnviron{guideblock}{%
  \ifdraft
    \par\medskip\noindent{\color{cGuide}\rule{\linewidth}{0.7pt}}\par\nobreak
    \begingroup\small\color{cGuide}\noindent\textbf{【写作提示】}\ignorespaces\BODY\par\endgroup
    \noindent{\color{cGuide}\rule{\linewidth}{0.7pt}}\par\medskip
  \fi
}

% 行内短注
\newcommand{\note}[1]{\ifdraft{\color{cNote}\small【注：#1】}\fi}

% 未算出的数字：一律用 \ph 占位，严禁编造数字
% 正文中允许在占位符内断行（避免长占位符把行撑出边界），数学模式中则用 \hbox 保持整体
\newcommand{\ph}[1][待计算]{%
  \ifmmode\hbox{\color{cPh}\bfseries 【#1】}%
  \else{\color{cPh}\bfseries 【#1】}\fi}

% 证据编号（证据账本 C-Qn-xx）
\newcommand{\evid}[1]{\ifdraft{\color{cEvid}(#1)}\else（见 #1）\fi}

% 口径编号（口径框 C-01…C-08）
\newcommand{\calref}[1]{口径 #1}

% 口径框（紧凑版：小字号、无首行缩进，正文用「定义/数据/时段/单位/不含/对照/复算」逐项写）
\NewEnviron{caliber}[2]{%
  \par\smallskip\noindent{\color{cEvid}\rule{\linewidth}{0.7pt}}\par\nobreak
  {\noindent\heiti 口径 #1\quad #2\par}\nobreak
  \begingroup\footnotesize\setlength{\parindent}{0pt}\BODY\par\endgroup
  \noindent{\color{cEvid}\rule{\linewidth}{0.7pt}}\par\smallskip
}

% 创新五段式
\NewEnviron{innovation}[2]{%
  \par\medskip\noindent{\heiti 创新 #1\quad #2}\par\nobreak\vspace{0.2ex}
  \begingroup\BODY\par\endgroup\medskip
}
\newcommand{\ip}[1]{\par\noindent\textbf{#1}\quad\ignorespaces}

% 插图：图片不存在时画红框占位，保证任何阶段都能编译
\newcommand{\includefig}[2][width=\linewidth]{%
  \IfFileExists{figures/#2}%
    {\includegraphics[#1]{figures/#2}}%
    {\fbox{\begin{minipage}[c][2.6cm][c]{0.9\linewidth}\centering\color{cPh}\small
      【缺图】\texttt{\detokenize{figures/#2}}\par\vspace{0.6ex}
      待由 \texttt{C题代码/code/plot\_fig*.py} 生成后放入 \texttt{figures/}\end{minipage}}}%
}

% 附录源程序：自动原样引入 code/ 目录下的真实文件，
% 保证"附录代码 = 真正跑出结果的代码"；文件缺失时给出红框提示而不中断编译。
\newcommand{\codefile}[2][py]{%
  \par\medskip\noindent{\heiti 源程序：\texttt{\detokenize{code/#2}}}\par\nobreak
  \IfFileExists{code/#2}%
    {\lstinputlisting[style=#1]{code/#2}}%
    {\begin{center}\fcolorbox{cPh}{white}{\parbox{0.92\linewidth}{\centering\color{cPh}\small
       【缺文件】\texttt{\detokenize{code/#2}}\par\vspace{0.4ex}
       请把源程序放入 \texttt{code/} 目录，附录会\textbf{自动原样引入}（不要用截图）}}\end{center}}%
  \par\medskip
}

% 论文标题（摘要页与 PDF 元数据共用，只改这一处）
\newcommand{\papertitle}{微网与外部电网电力调控策略}

% 单位（正体，与变量斜体区分）
\newcommand{\unit}[1]{\textup{#1}}

% 正文中引用题面表格
\newcommand{\tabref}[1]{表~#1}
